\documentclass[10pt,a4paper]{amsart}
\usepackage[margin=19mm]{geometry}
\usepackage{amsmath,amssymb,mathtools}
\usepackage[scr=rsfs,cal=euler]{mathalfa}
\usepackage{xfrac}
\usepackage[unicode,hidelinks,bookmarks=false]{hyperref}
\newcommand{\ie}{i.e.\ }
\newcommand{\cf}{cf.\ }
\newcommand{\resp}{respectively\ }
\newcommand{\Subsection}{\subsection}
\providecommand{\nobreakdash}{\nobreak\hbox{-}}
\setlength{\emergencystretch}{2em}
\multlinegap0pt
\usepackage{amssymb}
\usepackage{mathtools}
\usepackage{bm}
\usepackage{cancel}
\usepackage{float}
\usepackage{tikz}
\usepackage{xcolor}
\newtheorem{theorem}{Theorem}
\newtheorem{lemma}{Lemma}
\usepackage{cleveref}
\crefname{theorem}{Theorem}{Theorems}
\crefname{lemma}{Lemma}{Lemmas}
    \newcommand{\A}{\alg{A}}
    \newcommand{\B}{\alg{B}}
    \newcommand{\C}{\class{C}}
\newcommand{\Dia}{\Diamond}
\renewcommand{\Diamond}{\gooddiamond}
    \newcommand{\Kff}[2]{\logic{K4^{#1}_{#2}}}
    \renewcommand{\S}{\class{S}}
\newcommand{\Sub}{\mathsf{Sub}}
\newcommand{\alg}[1]{\mathfrak{#1}}
\newcommand{\class}[1]{\mathcal{#1}}
\newcommand{\emp}{\emptyset}
\newcommand{\gooddiamond}{\hspace{.2ex}\text{%
  \tikz[baseline=-.6ex, rounded corners=.01ex, rotate=45, line width=.12ex]
    {\draw(-.5ex,-.5ex) rectangle (.5ex,.5ex);}}\kern.2ex}
\newcommand{\inj}{\hookrightarrow}
\newcommand{\logic}[1]{\mathsf{#1}}
\renewcommand{\phi}{\varphi}
\title{Stable Canonical Rules and Formulas for \\ Pre-transitive Logics via Definable Filtration}
\author{Tenyo Takahashi}
\date{Source: arXiv:2510.22638v2 (CC BY 4.0)}
\begin{document}
\maketitle

\noindent\emph{Source and licence.} Stable Canonical Rules and Formulas for \\ Pre-transitive Logics via Definable Filtration. Version arXiv:2510.22638v2 (CC BY 4.0), distributed by arXiv under the Creative Commons Attribution 4.0 International licence, http://creativecommons.org/licenses/by/4.0/, as recorded in the arXiv licence metadata for this exact version and shown on the abstract page https://arxiv.org/abs/2510.22638v2. Full licence notice: \texttt{licenses/arxiv-2510.22638-CC-BY-4.0.txt}.\par\smallskip

\noindent\textit{Editorial scope.} Counted unit: the article's lemma pattern (source0.tex lines 730--737). The article's complete original proof of it is delivered with it (source0.tex lines 739--754, one \texttt{proof} environment of 16 lines). No mathematics is written, repaired or re-derived: the statement and the proof are the article's own lines, in the article's own notation, and no retained line is substituted or re-flowed.  Retained uncounted as the source-local context the proof needs: lines 556--559; lines 697--699; lines 724--726.  Nothing was omitted from any retained span, so the package carries no omission record.  No retained line contains a citation, so the wrapper carries no bibliography and no bibliography entry.  No retained line contains a figure environment, an image-inclusion command or drawing code, so no image is delivered and the package declares no asset dependency.  Editorial formatting only, in addition: an AMS \texttt{amsart} preamble that restates the article's own declarations and macros used by the retained lines, namely the environments \texttt{theorem}, \texttt{lemma} on separate counters, together with the standard packages those lines need.  The retained environments print in this document as Theorem 1, Lemma 1 in order of appearance, whereas the article prints the same items with the numbering of its own full document.  The complete native source of the article as distributed in the arXiv e-print for this exact version is bundled unchanged as \texttt{sources/187812/source0.tex}.
\begin{theorem} \label{3: Thm scr complete} 
    Let $\S$ be a rule system that admits definable filtration. For any rule $\rho$, there exist stable canonical rules $\rho(\A_1, D_1), \dots, \rho(\A_n, D_n)$ where each $\A_i$ is a finite $\S$-algebra and $D_i \subseteq A_i$, such that for any $\S$-algebra $\B$,
    \[\B \models \rho \text{ iff } \B \models\rho(\A_1, D_1), \dots, \rho(\A_n, D_n),\]
\end{theorem}

\begin{theorem} \label{3: Thm K4m1 admits filtration}
    For any $m \geq 1$, the logic $\Kff{m+1}{1}$ admits definable filtration. 
\end{theorem}

\begin{lemma} \label{3: Lem si reflect stable subalg}
    Let $\A$ be a finite modal algebra and $\B$ be a s.i.~modal algebra. If $\A$ is a stable subalgebra of $\B$, i.e., $\A \inj_\emp \B$, then $\A$ is also s.i.
\end{lemma}

\begin{lemma} \label{3: Lem K4m1 refutation pattern}
    For any formula $\phi$, there exist pairs $(\A_1, D_1), \dots, (\A_n, D_n)$ such that each $\A_i = (A_i, \Dia_i)$ is a finite s.i.~$\Kff{m+1}{1}$-algebra, $D_i \subseteq A_i$, and for any s.i.~modal algebra $\B = (B, \Dia)$, the following conditions are equivalent:
    \begin{enumerate}
        \item[(1)] $\B \not\models \phi$.
        \item[(2)] There is $1 \leq i \leq n$ and a stable embedding $h: \A_i \inj_{D_i} \B$.
        \item[(3)] There is a s.i.~homomorphic image $\C = (C, \Dia)$ of $\B$, $1 \leq i \leq n$, and a stable embedding $h: \A_i \inj_{D_i} \C$.
    \end{enumerate}
\end{lemma}

\begin{proof}
    If $\phi \in \Kff{m+1}{1}$, then let $n=0$. Assume that $\phi \notin \Kff{m+1}{1}$ and let $\Theta = \Sub(\phi)$. Since $\Kff{m+1}{1}$ admits definable filtration by \Cref{3: Thm K4m1 admits filtration}, there is a finite subformula-closed set $\Theta'$ containing $\Theta$ such that, for any $\Kff{m+1}{1}$-algebra $\B$ and any valuation $V$ on $\B$, there is a definable filtration $(\B', V')$ of $(\B, V)$ for $\Theta$ through $\Theta'$ such that $\B' \models \Kff{m+1}{1}$. Let $l = |\Theta'|$. Since Boolean algebras are locally finite, up to isomorphism, there are finitely many tuples $(\A, V, D)$ satisfying the following conditions:
    \begin{enumerate}
        \item $\A$ is a finite s.i.~$\Kff{m+1}{1}$-algebra based on an at most $l$-generated Boolean algebra and $\A \not\models \phi$,
        \item $V$ is a valuation on $\A$ such that $\A, V \not\models \phi$ and $V(p) = 0$ for $p \notin \Theta'$,
        \item $D = \{V(\psi): \Dia\psi \in \Theta\}$.
    \end{enumerate}
    Let $(\A_1, V_1, D_1), \dots, (\A_n, V_n, D_n)$ be an enumeration of such tuples. We show that $(\A_1, D_1), \dots, (\A_n, D_n)$ is the desired set of pairs. Let $\B$ be a s.i.~$\Kff{m+1}{1}$-algebra. 

    $(1) \Rightarrow (2)$. Suppose that $\B \not\models \phi$. Let $V$ be a valuation on $\B$ such that $\B, V \not\models \phi$. Then there is a definable filtration $(\B', V')$ of $(B, V)$ for $\Theta$ through $\Theta'$ such that $\B' \models \Kff{m+1}{1}$. By the definition of definable filtration, $\B'$ is a stable subalgebra of $\B$, so $\B'$ is also s.i.~by \Cref{3: Lem si reflect stable subalg}. Then the same argument as in the proof of \Cref{3: Thm scr complete} shows that the tuple $(\B', V', D)$ is identical to $(\A_i, V_i, D_i)$ for some $1 \leq i \leq n$. Since $\B' \inj_D \B$ by the definition of definable filtration, we conclude $\A_i \inj_{D_i} \B$.

    $(2) \Rightarrow (3)$. This is obvious by taking $\C = \B$.

    $(3) \Rightarrow (1)$. Suppose that there is a s.i.~homomorphic image $\C$ of $\B$, $1 \leq i \leq n$, and a stable embedding $h: \A_i \inj_{D_i} \C$. Let $V_i$ be valuation on $\A_i$ such that $\A_i, V_i \not\models \phi$. Define a valuation $V$ on $\C$ by $V(p) = h(V_i(p))$. The same argument as in the proof of \Cref{3: Thm scr complete} shows that $\C \not\models \phi$. Thus, $\B \not\models \phi$ since $\C$ is a homomorphic image of $\B$. 
    
\end{proof}

\end{document}
